\documentclass[preview]{standalone}

\usepackage{amsmath}
\usepackage{amssymb}
\usepackage{stellar}
\usepackage{definitions}

\begin{document}

\id{psicologia-forme-pensiero}
\genpage

\section{Forme del pensiero}

\begin{snippet}{forme-pensiero-introduzione}
    Si possono distinguere tre forme principali di pensiero: il
    \emph{pensiero verbale}, strutturato in forma linguistica; il
    \emph{pensiero per immagini}, basato su rappresentazioni analogiche; e il
    \emph{pensiero astratto}, formato da codici non riducibili a parole o
    immagini.
\end{snippet}

\begin{snippetdefinition}{linguaggio-subvocale-definition}{Linguaggio subvocale}
    Il \emph{linguaggio subvocale} è il linguaggio silenzioso che si attiva
    quando pensiamo come se parlassimo sottovoce a noi stessi.
\end{snippetdefinition}

\begin{snippet}{pensiero-verbale}
    Watson interpretava il pensiero come linguaggio subvocale, riducendolo a
    movimenti minimi dell'apparato fonatorio. Questa ipotesi fu indebolita da
    casi sperimentali in cui il pensiero risultava possibile anche in assenza di
    movimento muscolare volontario, mostrando che linguaggio parlato e pensiero
    possono essere dissociati.
\end{snippet}

\begin{snippet}{vygotskij-piaget-linguaggio}
    Vygotskij e Piaget studiarono il rapporto tra pensiero e linguaggio nello
    sviluppo. Per Vygotskij, il linguaggio infantile nasce da radici sociali: il
    linguaggio egocentrico deriva dal linguaggio sociale e conduce
    progressivamente al linguaggio interiore, centrale per il pensiero.
\end{snippet}

\begin{snippetdefinition}{immagine-mentale-definition}{Immagine mentale}
    Un'\emph{immagine mentale} è una rappresentazione psicologica di percezioni
    e azioni, generata dalla corrispondente mappa cerebrale momentanea di una
    data situazione.
\end{snippetdefinition}

\begin{snippetexample}{gatto-tappeto-example}{Il gatto sul tappeto}
    Se bisogna scegliere tra due rappresentazioni visive della frase ``il gatto
    è sul tappeto'', entrambe possono sembrare verbalmente corrette. Tuttavia,
    di solito si preferisce l'immagine in cui il gatto è orientato nel modo più
    abituale, perché corrisponde alla scena costruita spontaneamente nella
    mente.
\end{snippetexample}

\includesnpt[width=48\%,src=/snippet/static-psychology/mental-image-cat-on-mat.jpg]{centered-img}

\begin{snippetexample}{rotazione-mentale-example}{Rotazione mentale}
    In un classico compito di rotazione mentale, ai partecipanti vengono
    presentate lettere \texttt{R} normali o speculari ruotate tra \(0^\circ\) e
    \(180^\circ\). Il partecipante deve stabilire se la lettera sia normale o
    speculare. I tempi di reazione aumentano con l'angolo di rotazione,
    suggerendo che le persone ruotano mentalmente l'immagine fino alla posizione
    verticale prima di decidere.
\end{snippetexample}

\includesnpt[width=50\%,src=/snippet/static-psychology/mental-rotation-letter-r.jpg]{centered-img}

\begin{snippet}{percezione-immaginazione-visiva}
    Le immagini mentali vengono usate per rispondere a domande sul mondo e per
    costruire scene assenti. Studi di neuroimaging mostrano che aree cerebrali
    simili possono attivarsi durante la percezione visiva e durante la
    produzione di immagini mentali.
\end{snippet}

\includesnpt[width=62\%,src=/snippet/static-psychology/visual-perception-imagery-brain-areas.jpg]{centered-img}

\begin{snippetdefinition}{modello-mentale-spaziale-definition}{Modello mentale spaziale}
    Un \emph{modello mentale spaziale} è una rappresentazione interna della
    disposizione spaziale di oggetti, luoghi o eventi, costruita a partire dalla
    percezione o da descrizioni verbali.
\end{snippetdefinition}

\begin{snippetexample}{modello-mentale-spaziale-example}{Modello mentale spaziale}
    Quando leggiamo una descrizione di una stanza, possiamo immaginare la
    posizione degli oggetti al suo interno. Se ci proiettiamo mentalmente nella
    scena, rispondiamo più rapidamente a domande su ciò che sarebbe davanti a
    noi rispetto a ciò che sarebbe dietro.
\end{snippetexample}

\includesnpt[width=46\%,src=/snippet/static-psychology/spatial-mental-model.jpg]{centered-img}

\begin{snippetdefinition}{linguaggio-mente-definition}{Linguaggio della mente}
    Il \emph{linguaggio della mente} è, secondo Fodor, un sistema di
    rappresentazioni astratte e proposizionali che non coincide con parole o
    immagini e che costituirebbe la base del pensiero.
\end{snippetdefinition}

\begin{snippet}{caratteri-linguaggio-mente}
    Secondo Fodor, le rappresentazioni del linguaggio della mente hanno parti
    costituenti combinabili secondo una sintassi, sono composizionali e sono
    regolate da relazioni di implicazione. Questa ipotesi attribuisce al
    pensiero una forma astratta non derivata direttamente dalle modalità
    sensoriali.
\end{snippet}

\begin{snippetnote}{ipotesi-amodale-limite-note}{Limite dell'ipotesi amodale}
    A oggi non è emersa un'evidenza neuropsicologica chiara dell'esistenza di
    simboli amodali nel cervello.
\end{snippetnote}

\end{document}
